\documentclass[12pt]{article}
\usepackage[utf8]{inputenc}
\usepackage[T1]{fontenc}
%\usepackage{tamil}{}
\usepackage{geometry}
\usepackage{enumitem}
\usepackage{xcolor}
\usepackage{hyperref}
\usepackage{graphicx}
\usepackage{url}

% Set page margins
\geometry{a4paper, margin=1in}

% Define a custom color for headings
\definecolor{headcolor}{RGB}{139,69,19} % Saddle brown

% Customize section formatting
\usepackage{titlesec}
\titleformat{\section}
  {\normalfont\Large\bfseries\color{headcolor}}
  {\thesection}{1em}{}
\titleformat{\subsection}
  {\normalfont\large\bfseries\color{headcolor!80}}
  {\thesubsection}{1em}{}

% Disable fontspec as requested
\pagestyle{plain}

\begin{document}

\begin{titlepage}
    \centering
    \vspace*{2cm}
    {\Huge\bfseries\color{headcolor} Tamil Nadu's Enduring Legacy\par}
    \vspace{1cm}
    {\Large Culture, Manuscripts, and Language\par}
    \vspace{2cm}
    {\large A glimpse into a heritage often unseen\par}
    \vspace{3cm}
    \vfill
    {\today\par}
\end{titlepage}

\tableofcontents
\newpage

\section{Vibrant Culture}
Tamil Nadu, the land of the Tamils, possesses a cultural tapestry that is among the oldest continuously surviving classical civilizations in the world. While Indian popular culture often presents a homogenized view, the specific traditions of Tamil Nadu—from its temple architecture to its folk arts—remain deeply nuanced and regionally distinct.

\subsection{Classical Arts and Temple Traditions}
The culture is inextricably linked to its temple towns. The \textbf{Brihadeeswarar Temple} in Thanjavur, a UNESCO World Heritage Site, is not merely a monument but a living center of Chola-era traditions. The \textbf{Bharatanatyam} dance form, originating from the \textit{sadir} tradition of temple dancers (\textit{devadasis}), is a highly codified language of gesture and emotion, described in the ancient text \textit{ Natya Shastra} and the Tamil text \textit{Pancha Marabu}. Similarly, \textbf{Carnatic music}, with its 72 \textit{melakarta} ragas, finds its most ardent practitioners in Tamil Nadu. The \textit{Thiruppugazh} hymns of Arunagirinathar and the \textit{krithis} of the \textit{Tamil Trinity}—Tyagaraja, Muthuswami Dikshitar, and Syama Sastri—form the bedrock of this musical tradition, which is often transmitted through \textit{gurukula} systems in cities like Chennai (the "Detroit of Carnatic music") and Thanjavur \cite{krishnaswamy2001}.

\subsection{Folk Arts and Festivals}
Beyond the classical realm, a vibrant folk culture persists, often overlooked. \textbf{Therukoothu} (street theater) is a rural art form that dramatizes episodes from the \textit{Mahabharata} and \textit{Ramayana} with elaborate costumes, makeup, and a distinct blend of song and oratory. \textbf{Villu Paattu} (bow-song) is a narrative art where a bow-shaped string instrument is struck as a primary percussion, accompanied by a troupe of singers who weave stories of local deities and epics. Festivals like \textbf{Pongal}, the harvest festival, are celebrated with \textit{jallikattu} (the taming of bulls, a controversial but deeply rooted tradition) and \textit{kollam} (intricate floor art drawn with rice flour) that are distinctly Tamil, predating and differing from the northern \textit{rangoli} \cite{blackburn1988}.

\section{Manuscripts}
The manuscript heritage of Tamil Nadu is a vast, largely un-catalogued treasure that holds the keys to South Asian medicine, philosophy, and literature. These manuscripts, written on palm leaves (\textit{ola}) and later paper, represent a knowledge economy that operated for millennia outside the sphere of colonial and northern Indian historical narratives.

\subsection{Palm Leaf Manuscripts and the \textit{Sarasvati Mahal}}
The most renowned repository is the \textbf{Sarasvati Mahal Library} in Thanjavur, which holds over 60,000 palm leaf and paper manuscripts. Founded by the Nayak kings and enriched by the Maratha rulers, it contains a unique collection that includes the only known copies of several medieval medical texts (\textit{siddha vaidyam}), commentaries on the \textit{Tirukkural}, and rare illustrated manuscripts on music and architecture \cite{seshagiri2003}. The \textit{Siddha} medical system, one of India's oldest, is preserved almost exclusively in these palm leaf manuscripts, written in a cursive form of Tamil called \textit{Vatteluttu} and Grantha script. These texts detail complex alchemical preparations (\textit{rasavatham}) and treatments for chronic diseases—a heritage that remains scientifically under-explored due to linguistic and accessibility barriers.

\subsection{Preservation and Access}
A significant portion of Tamil manuscripts remains in private collections (\textit{maṭhas}) and family libraries across the Tanjore, Madurai, and Tirunelveli regions. The \textbf{Tamil Nadu Archives} in Chennai, one of the largest archival repositories in Asia, holds centuries of administrative records in Tamil, Persian, Marathi, and Telugu, chronicling the transition from the Vijayanagara Empire to the British Raj. However, as noted by scholars like David Shulman, the sheer volume and the fragility of these materials mean that the majority of this textual heritage is inaccessible to the general Indian public, with only a fraction digitized by institutions like the French Institute of Pondicherry \cite{shulman2016}.

\section{Local Language}
Tamil is not merely a language; it is a classical language with an unbroken literary tradition spanning over two millennia. However, its diglossic nature and the rich ecosystem of its dialects and sociolects are often hidden from the outside view.

\subsection{Classical Tamil and the \textit{Sangam} Literature}
The history of Tamil is divided into three periods: \textit{Sangam} (circa 500 BCE–300 CE), medieval, and modern. The \textbf{Sangam literature}, consisting of 2,381 poems attributed to over 473 poets, is the earliest extant secular literature in India. Texts like the \textit{Tolkāppiyam} (the earliest surviving Tamil grammar) codify not just linguistics but also ancient social and ecological concepts, classifying landscapes into five \textit{tinais} (thinai)—kurinji (mountains), mullai (forests), marutham (farmlands), neithal (coast), and paalai (desert)—each with its own deity, social structure, and poetic conventions. This literature, which predates the dominant Sanskritic \textit{itihasa} tradition, represents a parallel classical world that is foundational to Tamil identity \cite{zvelebil1973}.

\subsection{Diglossia, Dialects, and Script Evolution}
A unique feature of Tamil is its strong diglossia: the formal, literary variety used in media and education differs significantly from the colloquial varieties spoken in everyday life. There are distinct regional dialects—such as those of Kongu (western region), Madurai, Tirunelveli, and the Brahmin and non-Brahmin sociolects—that encode a vast cultural and social history. The Tamil script itself evolved from the ancient \textbf{Grantha} and \textbf{Vatteluttu} scripts, with the modern script being standardized in the 20th century. Unlike most Indian languages, Tamil has a strong purist movement that has resisted the influx of Sanskrit and English loanwords, preserving a unique linguistic identity. The language's classical status, granted by the Government of India in 2004, officially recognized its antiquity and independence as a literary culture, yet the depth of its dialectal and manuscript diversity remains hidden to the majority of Indians who are only familiar with its standardized, formal register \cite{annamalai2011}.

\begin{thebibliography}{9}
    \bibitem{krishnaswamy2001}
    Krishnaswamy, S. (2001). \textit{The Tamil Trinity of Carnatic Music}. Madras: Institute of Asian Studies.

    \bibitem{blackburn1988}
    Blackburn, S. H. (1988). \textit{Singing of Birth and Death: Texts in Performance}. Philadelphia: University of Pennsylvania Press. (On Villu Paattu and Therukoothu).

    \bibitem{seshagiri2003}
    Seshagiri, S. (2003). \textit{The Sarasvati Mahal Library: A History}. Thanjavur: Thanjavur Maharaja Serfoji's Sarasvati Mahal Library.

    \bibitem{shulman2016}
    Shulman, D. (2016). \textit{Tamil: A Biography}. Cambridge, MA: Harvard University Press. (Insights on manuscript culture and the nature of Tamil).

    \bibitem{zvelebil1973}
    Zvelebil, K. V. (1973). \textit{The Smile of Murugan: On Tamil Literature of South India}. Leiden: E. J. Brill. (Comprehensive analysis of Sangam literature).

    \bibitem{annamalai2011}
    Annamalai, E. (2011). \textit{Social Dimensions of Modern Tamil}. Pondicherry: French Institute of Pondicherry. (On diglossia, dialects, and language policy).

\end{thebibliography}

\end{document}